% TOP_PROBLEM: {"id": 30003641, "title": "Partition Regularity of Pythagorean Triples", "category_id": 1, "category": {"id": 1, "name": "number_theory", "display_name": "Number Theory", "description": "Properties of integers, prime numbers, Diophantine equations.", "slug": "number-theory", "order_index": 1, "created_at": "2026-07-31T15:26:25.670Z"}, "status": "open"}
% FIRST_POSTED: 2026-09-25
% ENTRY_KIND: statement_only
% !TeX program = lualatex
% Definition and sources only; this file is not an LLM solution attempt.
\documentclass[11pt]{article}
\usepackage[margin=25mm]{geometry}
\usepackage{fontspec}
\setmainfont{Latin Modern Roman}
\usepackage{amsmath,amssymb,mathtools}
\usepackage{unicode-math}
\setmathfont{Latin Modern Math}
\usepackage{xurl,hyperref}
\hypersetup{colorlinks=true,urlcolor=blue}
\setcounter{secnumdepth}{0}
\setlength{\parindent}{0pt}
\setlength{\parskip}{5pt}
\setlength{\emergencystretch}{3em}
\begin{document}
\section*{\texorpdfstring{Partition Regularity of Pythagorean Triples}{Partition Regularity of Pythagorean Triples}}\label{30003641}
\subsection{Definitions and mathematical statement}
A finite colouring of the positive integers is a function $c:{\mathbb{Z}}_{\ge1}\to\{1,\ldots,r\}$ for some finite $r\ge1$. The Pythagorean equation is partition regular if
\[
 \forall r\ \forall c\ \exists x,y,z\in{\mathbb{Z}}_{\ge1},
 \qquad x^2+y^2=z^2,
 \qquad c(x)=c(y)=c(z).
\]
Determine whether this assertion holds. The triple need not be primitive, so a common divisor of its entries is allowed. Positivity excludes zero solutions. The number of colours is arbitrary and fixed for each colouring, not just two; a theorem for one specific number of colours does not settle the universal finite-colouring statement.
\par\smallskip\textit{Formulation and concept references:} \href{https://arxiv.org/abs/2411.17523}{[S1]}.

\subsection{Short English statement}
Does every finite colouring of the positive integers contain a single-coloured Pythagorean triple?
\subsection{Sources}
\begin{itemize}
\item[S1]Partition regularity of homogeneous quadratics: Current trends and challenges. Submitted scholarly source, accessed 2026-08-31.
\url{https://arxiv.org/abs/2411.17523}.
\textit{Formulation source.}
\end{itemize}
\end{document}
